% Source: 05 - Tue Oct 6 2026.pdf, page 2. Content remains in source order.
\sourcepage{05}{2}
For each gate $g\in G$ we introduce a sub-formula $\Gamma_g$ describing the input-output relation of the gate:
\begin{itemize}
\item AND gate:
\sourcefigure[0.3]{assets/05/page-002-and-gate.png}
\[
\Gamma_g:\ y_g\underbrace{\leftrightarrow}_{\text{``equation''}}x_{g_1}\land\cdots\land x_{g_k}\land y_{g_1}\land\cdots\land y_{g_h}.
\]
\item OR gate:
\sourcefigure[0.3]{assets/05/page-002-or-gate.png}
\[
\Gamma_g:\ y_g\leftrightarrow x_{g_1}\lor\cdots\lor x_{g_k}\lor y_{g_1}\lor\cdots\lor y_{g_h}.
\]
\item NOT gate:
\sourcefigure[0.3]{assets/05/page-002-not-gate.png}
\[\Gamma_g:\ y_g\leftrightarrow\neg(z_{g_1}).\]
\end{itemize}
Observe that subexpression $\Gamma_g$ is true under a truth assignment iff $y_g$ is assigned the correct output value w.r.t. the values assigned to the gate's inputs.

We can then obtain the formula
\[\psi(\vec x,\vec y)=\bigwedge_{g\in G}\Gamma_g.\]
Formula $\psi$ describes a configuration of the circuit $C$:
\[
\begin{aligned}
\psi(\vec b_1,\vec b_2)=1
&\Longleftrightarrow\text{the values of the outputs of the gates of }C(\vec b_1)\text{ are }(b_2)_g:g\in G\\
&\Longleftrightarrow(\vec b_1,\vec b_2)\text{ is a valid configuration.}
\end{aligned}
\]
